\hypertarget{appendix-1bc40b08db4a}{}\label{referencestart:appendix-1bc40b08db4a}
\section{Hardware depth-buffer linearization}\label{appendix-1bc40b08db4a-hardware-depth-buffer-linearization}

\hypertarget{appendix-1bc40b08db4a-hardware-depth-buffer-linearization}{}

\textbf{Curriculum depth:} Developed · \textbf{First seen:} Phase 11 · \textbf{Latest development:} Phase 11

\begin{allmightroledefinition}

\subsection{Definition and mechanism}\label{appendix-1bc40b08db4a-definition-and-mechanism}

\hypertarget{appendix-1bc40b08db4a-definition-and-mechanism}{}

Perspective projection maps axial distance nonlinearly into NDC depth. For a declared projection depth row \(q=-A+B/d\), solve \(d=B/(q+A)\); finite conventional, finite reversed, infinite reversed, and OpenGL-style conventions have different coefficients or NDC conversions. The depth sample must first be interpreted in the same range and direction used by the projection matrix. This is why a single generic formula without an API and projection contract is unsafe. Verify near and far endpoints, test round trips over a range of distances, and handle an infinite-far endpoint such as reversed depth \(q=0\) explicitly. Near/far endpoint checks catch sign mistakes immediately; interior round trips then test scaling and finite precision.

\end{allmightroledefinition}

\subsection{Prerequisites}\label{appendix-1bc40b08db4a-prerequisites}

\hypertarget{appendix-1bc40b08db4a-prerequisites}{}

\begin{itemize}
\tightlist
\item
  \href{../04-Light-Transport-and-Image-Synthesis/main.pdf\#nameddest=appendix-1c07d64f5e70}{Camera view and clip-space pipeline (Light Transport and Image Synthesis, Reference notes: Camera view and clip-space pipeline)}
\item
  \hyperref[appendix-cac01ebeeb11]{Axial depth and Euclidean range}
\item
  \href{../04-Light-Transport-and-Image-Synthesis/main.pdf\#nameddest=appendix-97578ed97952}{Depth buffering and visibility order (Light Transport and Image Synthesis, Reference notes: Depth buffering and visibility order)}
\end{itemize}

\subsection{Used by}\label{appendix-1bc40b08db4a-used-by}

\hypertarget{appendix-1bc40b08db4a-used-by}{}

\begin{itemize}
\tightlist
\item
  \hyperref[appendix-2cbf2968e866]{Position reconstruction from depth}
\end{itemize}

\subsection{Full teaching explanation}\label{appendix-1bc40b08db4a-full-teaching-explanation}

\hypertarget{appendix-1bc40b08db4a-full-teaching-explanation}{}

\begin{itemize}
\tightlist
\item
  \href{../03-Geometry-Cameras-and-Image-Formation/main.pdf\#nameddest=02-foundations-ii-ch035-recover-geometry-from-depth-and-construct-screen-space-rays}{Recover geometry from depth and construct screen-space rays (Geometry, Cameras, and Image Formation, Ch. 4)}
\end{itemize}

\subsection{Source coverage}\label{appendix-1bc40b08db4a-source-coverage}

\hypertarget{appendix-1bc40b08db4a-source-coverage}{}

Supplied lesson parts: 1079.

\subsection{Future keywords}\label{appendix-1bc40b08db4a-future-keywords}

\hypertarget{appendix-1bc40b08db4a-future-keywords}{}

standard depth, reversed-Z, infinite far plane
